\chapter{Phân tích và thiết kế hệ thống}
\label{ch:thiet-ke}

\section{Mô tả bài toán}

Sinh viên SGU hỏi bằng tiếng Việt --- thường viết tắt, sai chính tả, hỏi nối tiếp --- về quy chế đào tạo,
và nhận câu trả lời ngắn kèm văn bản và Điều nguồn để tự kiểm tra. Kho tài liệu gồm 6 văn bản đang áp
dụng, trong đó có văn bản sửa đổi (QĐ 3183/2025) chỉ áp dụng một phần cho khóa tuyển sinh từ 2025.

\section{User story}

Các câu chuyện người dùng dưới đây lấy từ câu hỏi thật trong nhật ký hỏi đáp:

\begin{itemize}
  \item Là sinh viên, tôi muốn hỏi ``thời gian tối đa hoàn thành khóa học'' và nhận câu trả lời kèm Điều để tự kiểm tra.
  \item Là sinh viên, khi hỏi điều không có trong quy chế (``bao nhiêu tiền một tín chỉ''), tôi muốn hệ thống nói rõ là không tìm thấy.
  \item Là sinh viên khóa 2025, khi hỏi ``8,4 hệ 10 thì hệ 4 được bao nhiêu'', tôi muốn câu trả lời theo thang điểm áp dụng cho khóa của tôi.
  \item Là sinh viên, tôi muốn hỏi bằng câu không dấu hoặc sai chính tả (``xuất sắt'', ``dạt diều khiện'') mà vẫn được trả lời.
  \item Là sinh viên, sau khi hỏi về cử nhân, tôi muốn hỏi tiếp ``còn kỹ sư?'' mà không phải lặp lại cả câu.
\end{itemize}

Hai story cuối mới được đáp ứng một phần (BM25 bỏ dấu) hoặc chưa (không nhớ hội thoại).

\section{Thiết kế}

\subsection{Sơ đồ pipeline}

Pipeline thật gồm sáu giai đoạn (Hình~\ref{fig:pipeline}, Chương~\ref{ch:moi-truong}): dữ liệu
$\rightarrow$ Node $\rightarrow$ embedding/index $\rightarrow$ truy hồi $\rightarrow$ tổng hợp (ngữ cảnh + LLM
+ trích nguồn) $\rightarrow$ đánh giá.

\subsection{Mô hình dữ liệu}

\begin{table}[H]
  \centering
  \small
  \begin{tabularx}{\textwidth}{l X c c}
    \toprule
    \textbf{Khóa metadata} & \textbf{Ý nghĩa} & \textbf{Embed} & \textbf{Prompt} \\
    \midrule
    \texttt{nguon\_trich} & ``Điều 9 --- Quy chế \ldots 2021'' (sinh tự động) & có & có \\
    \texttt{van\_ban} & Tên hiển thị của văn bản (\texttt{src/corpus.py}) & không & có \\
    \texttt{chuong} & Tên Chương & không & có \\
    \texttt{dieu} / \texttt{sua\_doi\_dieu} & Số Điều, hoặc số Điều được sửa đổi & không & có \\
    \texttt{phan} & \texttt{chinh} / \texttt{sua\_doi} / \texttt{ban\_hanh} / \texttt{noi\_dung} & không & không \\
    \texttt{phan\_doan} & Thứ tự đoạn khi một Điều bị cắt tiếp & không & có \\
    \texttt{file}, \texttt{file\_goc}, \texttt{url} & Truy vết về file tầng 2 và PDF gốc & không & không \\
    \bottomrule
  \end{tabularx}
  \caption{Metadata của Node (chỉ khóa có giá trị mới được gắn).}
\end{table}

Bộ câu hỏi truy hồi là JSON \texttt{\{"question", "dieu\_dung"\}}; bộ chấm câu trả lời thêm
\texttt{pham\_vi} (``trong''/``ngoai''), \texttt{dap\_an} và \texttt{nguon}.

\subsection{Luồng truy vấn}

\begin{figure}[H]
  \centering
  \begin{tikzpicture}[
      node distance=4mm,
      box/.style={rectangle, draw, rounded corners, minimum height=1cm, align=center, font=\small, text width=1.9cm},
      opt/.style={box, dashed},
      arr/.style={-{Stealth}, thick}
    ]
    \node[box] (q) {Câu hỏi};
    \node[box, right=of q] (v) {Vector top-10};
    \node[opt, right=of v] (b) {+ BM25\\ (RRF)};
    \node[opt, right=of b] (r) {Rerank\\ top-3};
    \node[box, right=of r] (p) {Prompt\\ QA\_PROMPT};
    \node[box, below=8mm of p] (l) {Ollama\\ qwen3};
    \node[box, left=of l] (a) {Trả lời\\ + Nguồn};
    \draw[arr] (q) -- (v); \draw[arr] (v) -- (b); \draw[arr] (b) -- (r); \draw[arr] (r) -- (p);
    \draw[arr] (p) -- (l); \draw[arr] (l) -- (a);
  \end{tikzpicture}
  \caption{Luồng truy vấn; khung nét đứt là bước tuỳ chọn, mặc định tắt.}
  \label{fig:luong-truy-van}
\end{figure}
